\section{Global maximizers for the Faber--Krahn problem in the Gaussian STFT setting}\label{sec:global-max}

In this section we prove existence of global maximizers for the Gaussian
time-frequency concentration problem under a measure constraint, and then
combine that existence result with the local rigidity theorem from the previous
section to recover the disk as the unique maximizer, up to translation. The
strategy is by now classical in time-frequency analysis: a profile decomposition
in the spirit of P.-L.~Lions' concentration-compactness, transplanted to the
Fock space via the Bargmann transform, dichotomizes the loss of compactness of a
maximizing sequence into the action of the Weyl translation group. A vanishing
remainder is then ruled out by an iterative extraction whose maximality forces
all the energy to remain captured by the extracted profiles. Closely related
arguments have appeared in the time-frequency setting in
\cite{Nicola-Romero-Trapasso, Marceca-Romero-Speckbacher, Fulsche-Hagger,
Samuelsen}, and the disk-rigidity result we recover is the celebrated theorem
of Nicola--Tilli \cite{Nicola-Tilli, Nicola-Tilli-2}.

Since the Bargmann transform is unitary, the variational problem may be written
on the Fock space side as
\[
M(s):=\sup_{|\Omega|=s}\lambda_{1,\varphi}(\Omega)
=
\sup_{\|F\|_{\mathcal F^2}=1} I_F(s).
\]
Here, and throughout the section, we exploit the translation invariance
\eqref{eq:I-translation-global} together with the pointwise identity
\eqref{eq:u-translation-global} (cf.\ also \eqref{eq:uF-translate}), the
basic regularity properties recorded in Lemma~\ref{lem:basic-properties-global},
and the elementary properties of the Fock space gathered in
Lemma~\ref{lem:Fock-basic}.

We begin with a Weyl-orthogonality lemma that captures the asymptotic
decoupling of two functions whose translation parameters drift apart. This
plays the role of \emph{Pythagorean orthogonality at infinity} on the Fock
space.

\begin{lemma}\label{lem:weyl-orthogonality-global}
Let \(F,G\in \mathcal F^2(\C)\), and let \(\{a_k\}\subset \C\) satisfy
\(|a_k|\to \infty\). Then
\[
T_{a_k}G \rightharpoonup 0
\qquad\text{weakly in }\mathcal F^2(\C),
\]
and in particular
\[
\langle F,T_{a_k}G\rangle_{\mathcal F^2}\to 0.
\]
Consequently,
\[
\|F+T_{a_k}G\|_{\mathcal F^2}^2
=
\|F\|_{\mathcal F^2}^2+\|G\|_{\mathcal F^2}^2+o(1).
\]
More generally, if \(|a_k-b_k|\to\infty\), then
\[
\langle T_{a_k}F,T_{b_k}G\rangle_{\mathcal F^2}\to 0.
\]
\end{lemma}

\begin{proof}
Using unitarity of the Weyl operators \(T_a\) on \(\mathcal F^2(\C)\) (cf.\
\eqref{eq:Ta}),
\[
\langle F,T_{a_k}G\rangle_{\mathcal F^2}
=
\langle T_{-a_k}F,G\rangle_{\mathcal F^2}.
\]
Now \(T_{-a_k}F\) is bounded in \(\mathcal F^2\) by unitarity, and the
identity \eqref{eq:u-translation-global} (equivalently
\eqref{eq:uF-translate}) yields
\[
u_{T_{-a_k}F}(z)=u_F(z+a_k)\qquad (z\in\C),
\]
so that the densities \(u_{T_{-a_k}F}\) drift uniformly to infinity.
Equivalently, since \(u_F\in C_0(\C)\) by Lemma~\ref{lem:Fock-basic}(iii), for
every compact \(K\subset \C\),
\[
\sup_{z\in K}|T_{-a_k}F(z)|e^{-\pi |z|^2/2}
=
\sup_{z\in K}u_{T_{-a_k}F}(z)^{1/2}
=
\sup_{z\in K}u_F(z+a_k)^{1/2}\to 0.
\]
Hence \(T_{-a_k}F\to 0\) locally uniformly. Combining this with the
boundedness of \(\{T_{-a_k}F\}\) in \(\mathcal F^2(\C)\), and using the fact
that the linear span of the reproducing kernels \(\{K_w\}_{w\in\C}\) is dense
in \(\mathcal F^2\), we conclude that \(T_{-a_k}F\rightharpoonup 0\) weakly in
\(\mathcal F^2\); in particular \(\langle T_{-a_k}F,G\rangle\to 0\). This
proves the first claim.

For the norm decoupling, expand the square:
\[
\|F+T_{a_k}G\|_{\mathcal F^2}^2
=
\|F\|_{\mathcal F^2}^2+\|T_{a_k}G\|_{\mathcal F^2}^2
+2\Re\langle F,T_{a_k}G\rangle_{\mathcal F^2}
=
\|F\|_{\mathcal F^2}^2+\|G\|_{\mathcal F^2}^2+o(1),
\]
where we used unitarity of \(T_{a_k}\) and the cross-term decay just proved.

Finally, the general case reduces to the first by writing
\[
\langle T_{a_k}F,T_{b_k}G\rangle_{\mathcal F^2}
=
\langle F,T_{-a_k}T_{b_k}G\rangle_{\mathcal F^2}
=
e^{i\theta_k}\langle F,T_{b_k-a_k}G\rangle_{\mathcal F^2},
\]
for some real phase \(\theta_k\) coming from the cocycle of the Weyl
representation; since this phase factor has modulus one, the conclusion
follows from \(|b_k-a_k|\to\infty\).
\end{proof}

The next lemma compares the density of a finite sum of widely separated
translates with the density of each individual translate. The precise statement
is what makes the profile decomposition compatible with the variational
quantities \(I_F(s)\), since these depend only on the densities and not on
the phases of the underlying functions.

\begin{lemma}\label{lem:pointwise-decoupling-global}
Fix \(J\ge 1\), let \(G^{(1)},\dots,G^{(J)}\in \mathcal F^2(\C)\), and let
\(\{a_k^{(j)}\}_{k\ge 1}\subset \C\) for \(1\le j\le J\) satisfy
\[
|a_k^{(j)}-a_k^{(\ell)}|\to\infty
\qquad\text{whenever }j\neq \ell.
\]
For each \(k\), set
\[
S_k:=\sum_{j=1}^J T_{a_k^{(j)}}G^{(j)}.
\]
Then the following hold.
\begin{enumerate}
\item[(i)] For every \(R>0\) and every \(1\le j\le J\),
\[
\sup_{z\in D_R(a_k^{(j)})}
\big|u_{S_k}(z)-u_{T_{a_k^{(j)}}G^{(j)}}(z)\big|
\to 0
\qquad (k\to\infty).
\]

\item[(ii)] For every \(R>0\), if
\[
E_{k,R}:=\C\setminus \bigcup_{j=1}^J D_R(a_k^{(j)}),
\]
then
\[
\limsup_{k\to\infty}\sup_{z\in E_{k,R}}u_{S_k}(z)
\le
\left(
\sum_{j=1}^J
\sup_{\zeta\in \C\setminus D_R} u_{G^{(j)}}(\zeta)^{1/2}
\right)^2.
\]
In particular, given \(\varepsilon>0\), there exists \(R_\varepsilon>0\) such that
for every \(R\ge R_\varepsilon\),
\[
\sup_{z\in E_{k,R}}u_{S_k}(z)\le \varepsilon
\]
for all sufficiently large \(k\).
\end{enumerate}
\end{lemma}

\begin{proof}
Throughout the proof we use repeatedly the identity
\[
u_{T_a F}(z)=u_F(z-a),
\qquad z,a\in \C,
\]
which is \eqref{eq:u-translation-global} (equivalently
\eqref{eq:uF-translate}); together with the elementary subadditivity
\[
u_{F+G}(z)^{1/2}\le u_F(z)^{1/2}+u_G(z)^{1/2},
\qquad z\in\C,
\]
which follows from the definition \(u_F=|F|^2 e^{-\pi |z|^2}\) and the
triangle inequality applied to \(|F+G|\le |F|+|G|\).

\medskip
\noindent\emph{Proof of (i).}
Fix \(R>0\) and \(j\in\{1,\dots,J\}\). Write
\[
S_k=T_{a_k^{(j)}}G^{(j)}+H_{k,j},
\qquad
H_{k,j}:=\sum_{\ell\neq j} T_{a_k^{(\ell)}}G^{(\ell)}.
\]
For every \(z\in \C\), the elementary inequality
\(\big||a+b|^2-|a|^2\big|\le 2|a|\,|b|+|b|^2\) yields
\begin{align*}
\big|u_{S_k}(z)-u_{T_{a_k^{(j)}}G^{(j)}}(z)\big|
&=
e^{-\pi |z|^2}
\big||T_{a_k^{(j)}}G^{(j)}(z)+H_{k,j}(z)|^2
-|T_{a_k^{(j)}}G^{(j)}(z)|^2\big| \\
&\le
2\,u_{T_{a_k^{(j)}}G^{(j)}}(z)^{1/2}u_{H_{k,j}}(z)^{1/2}
+u_{H_{k,j}}(z).
\end{align*}
By Lemma~\ref{lem:Fock-basic}(ii),
\[
u_{T_{a_k^{(j)}}G^{(j)}}(z)\le \|G^{(j)}\|_{\mathcal F^2}^2
\qquad \text{for every }z\in\C\text{ and every }k.
\]
Thus, to prove (i) it is enough to show that
\[
\sup_{z\in D_R(a_k^{(j)})} u_{H_{k,j}}(z)\to 0.
\]

For \(z\in D_R(a_k^{(j)})\) and \(\ell\neq j\),
\[
z-a_k^{(\ell)}\in D_R(a_k^{(j)}-a_k^{(\ell)}).
\]
Since \(|a_k^{(j)}-a_k^{(\ell)}|\to\infty\),
\[
\operatorname{dist}\big(D_R(a_k^{(j)}-a_k^{(\ell)}),0\big)
\ge |a_k^{(j)}-a_k^{(\ell)}|-R\to\infty.
\]
Because \(u_{G^{(\ell)}}\in C_0(\C)\) by Lemma~\ref{lem:Fock-basic}(iii),
\[
\sup_{z\in D_R(a_k^{(j)})} u_{T_{a_k^{(\ell)}}G^{(\ell)}}(z)
=
\sup_{\zeta\in D_R(a_k^{(j)}-a_k^{(\ell)})} u_{G^{(\ell)}}(\zeta)
\xrightarrow[k\to\infty]{} 0.
\]
The subadditivity of \(u_F^{1/2}\) yields
\[
u_{H_{k,j}}(z)^{1/2}
\le
\sum_{\ell\neq j} u_{T_{a_k^{(\ell)}}G^{(\ell)}}(z)^{1/2},
\]
so the preceding convergence implies
\[
\sup_{z\in D_R(a_k^{(j)})} u_{H_{k,j}}(z)^{1/2}\to 0.
\]
This proves (i).

\medskip
\noindent\emph{Proof of (ii).}
Fix \(R>0\) and \(z\in E_{k,R}\). By definition, \(z\notin D_R(a_k^{(j)})\)
for every \(j\), and hence \(|z-a_k^{(j)}|\ge R\). Therefore
\[
u_{T_{a_k^{(j)}}G^{(j)}}(z)
=
u_{G^{(j)}}(z-a_k^{(j)})
\le
\sup_{\zeta\in \C\setminus D_R}u_{G^{(j)}}(\zeta).
\]
Using again the subadditivity of \(u_F^{1/2}\),
\[
u_{S_k}(z)^{1/2}
\le
\sum_{j=1}^J u_{T_{a_k^{(j)}}G^{(j)}}(z)^{1/2}
\le
\sum_{j=1}^J \sup_{\zeta\in \C\setminus D_R}u_{G^{(j)}}(\zeta)^{1/2}.
\]
Squaring yields
\[
u_{S_k}(z)
\le
\left(
\sum_{j=1}^J \sup_{\zeta\in \C\setminus D_R}u_{G^{(j)}}(\zeta)^{1/2}
\right)^2.
\]
Taking the supremum over \(z\in E_{k,R}\) proves the displayed estimate in
(ii). Since each \(u_{G^{(j)}}\in C_0(\C)\), the right-hand side tends to
\(0\) as \(R\to\infty\); the final claim follows by choosing
\(R_\varepsilon\) so that the displayed bracket is at most
\(\sqrt{\varepsilon}\).
\end{proof}

The next lemma is a robust uniform-density perturbation estimate. Concretely,
it says that adding to a bounded sequence in \(\mathcal F^2(\C)\) a residue
whose density vanishes uniformly leaves the density essentially unchanged in
\(L^\infty\). This will be applied with \(S_k\) the truncated profile sum and
\(R_k\) the residual of the profile decomposition.

\begin{lemma}\label{lem:density-perturbation-global}
Let \(\{S_k\}\subset \mathcal F^2(\C)\) be bounded and let
\(\{R_k\}\subset \mathcal F^2(\C)\) satisfy
\[
\|u_{R_k}\|_{L^\infty(\C)}\to 0.
\]
Then
\[
\|u_{S_k+R_k}-u_{S_k}\|_{L^\infty(\C)}\to 0.
\]
More precisely, if \(\sup_k \|S_k\|_{\mathcal F^2}\le M\), then
\[
\|u_{S_k+R_k}-u_{S_k}\|_{L^\infty(\C)}
\le
2M\,\|u_{R_k}\|_{L^\infty(\C)}^{1/2}
+\|u_{R_k}\|_{L^\infty(\C)}.
\]
\end{lemma}

\begin{proof}
For every \(z\in \C\),
\begin{align*}
\big|u_{S_k+R_k}(z)-u_{S_k}(z)\big|
&=
e^{-\pi |z|^2}\big||S_k(z)+R_k(z)|^2-|S_k(z)|^2\big| \\
&\le
2\,u_{S_k}(z)^{1/2}u_{R_k}(z)^{1/2}
+u_{R_k}(z),
\end{align*}
again by the elementary inequality
\(\big||a+b|^2-|a|^2\big|\le 2|a|\,|b|+|b|^2\). By
Lemma~\ref{lem:Fock-basic}(ii),
\[
u_{S_k}(z)\le \|S_k\|_{\mathcal F^2}^2\le M^2
\qquad (z\in\C).
\]
Hence
\[
\big|u_{S_k+R_k}(z)-u_{S_k}(z)\big|
\le
2M\,\|u_{R_k}\|_{L^\infty}^{1/2}
+\|u_{R_k}\|_{L^\infty}.
\]
Taking the supremum over \(z\) yields the claim.
\end{proof}

The next lemma supplies the soft input on the function \(s\mapsto I_F(s)\)
that will be used at the end of the variational argument. It is essentially a
consequence of the bathtub principle and the basic regularity of \(u_F\),
already noted in Lemma~\ref{lem:basic-properties-global}.

\begin{lemma}\label{lem:I-continuity-global}
Let \(F\in \mathcal F^2(\C)\). Then the function
\[
s\mapsto I_F(s)
\]
is finite, nondecreasing, and continuous on \([0,\infty)\).
\end{lemma}

\begin{proof}
By Lemma~\ref{lem:Fock-basic}(ii)--(iii) together with the identity
\(\int_\C u_F\,dA=\|F\|_{\mathcal F^2}^2\), we have
\(u_F\in C_0(\C)\cap L^1(\C)\), so \(I_F(s)\) is finite for every \(s\ge 0\)
and bounded above by \(\|F\|_{\mathcal F^2}^2\). Monotonicity is immediate
from the definition of \(I_F(s)\) as a supremum over \(|\Omega|=s\): if
\(s_1\le s_2\), every admissible set for \(s_1\) can be enlarged to a set of
measure \(s_2\) without decreasing the integral, since \(u_F\ge 0\).

To prove continuity, write \(u:=u_F\) and let \(u^*\) denote the decreasing
rearrangement of \(u\) on \([0,\infty)\). Since \(u\in C_0(\C)\cap L^1(\C)\),
\(u^*\in L^1([0,\infty))\) and is equimeasurable with \(u\). The bathtub
principle yields
\[
I_F(s)=\int_0^s u^*(\tau)\,d\tau,\qquad s\ge 0.
\]
The right-hand side is the integral of an
\(L^1_{\mathrm{loc}}([0,\infty))\) function over \([0,s]\) and is therefore
absolutely continuous, hence continuous, in \(s\).
\end{proof}

The next result is the profile decomposition adapted to the present setting,
in the spirit of P.-L.~Lions' concentration-compactness method. Although the
overall scheme is by now standard, we spell out the points used in the
sequel; the version stated here is a Fock-space variant of profile
decompositions used in time-frequency analysis, see e.g.\
\cite{Nicola-Romero-Trapasso, Marceca-Romero-Speckbacher, Fulsche-Hagger,
Samuelsen}.

\begin{proposition}[Profile decomposition in the Fock space]\label{prop:profile-global}
Let \(\{F_k\}_{k\ge 1}\subset \mathcal F^2(\C)\) be a bounded sequence. Then,
after passing to a subsequence, there exist profiles
\[
G^{(j)}\in \mathcal F^2(\C),
\qquad j\ge 1,
\]
and sequences of points
\[
a_k^{(j)}\in \C,
\qquad j\ge 1,\ k\ge 1,
\]
such that:
\begin{enumerate}
    \item[(i)] for each \(j\neq \ell\),
    \[
    |a_k^{(j)}-a_k^{(\ell)}|\to \infty
    \qquad (k\to \infty);
    \]
    \item[(ii)] for every \(J\ge 1\), if
    \[
    R_k^{(J)}
    :=
    F_k-\sum_{j=1}^J T_{a_k^{(j)}}G^{(j)},
    \]
    then
    \[
    T_{-a_k^{(j)}}R_k^{(J)}\to 0
    \qquad\text{locally uniformly on }\C
    \]
    for every \(1\le j\le J\);
    \item[(iii)] for every \(J\ge 1\),
    \[
    \|F_k\|_{\mathcal F^2}^2
    =
    \sum_{j=1}^J \|G^{(j)}\|_{\mathcal F^2}^2
    +\|R_k^{(J)}\|_{\mathcal F^2}^2
    +o_k(1);
    \]
    \item[(iv)]
    \[
    \lim_{J\to\infty}\left(\limsup_{k\to\infty}
    \|u_{R_k^{(J)}}\|_{L^\infty(\C)}\right)=0.
    \]
\end{enumerate}
\end{proposition}

\begin{proof}
We argue by the standard iterative extraction scheme, but spell out the points
used in the rest of the section. Throughout, we say that a function in
\(\mathcal F^2(\C)\) is \emph{nontrivial} if it is not identically zero.

\medskip
\noindent\emph{Extraction of the first profile.}
If
\[
\limsup_{k\to\infty}\|u_{F_k}\|_{L^\infty(\C)}=0,
\]
there is nothing to prove: take all profiles equal to zero, all centers
arbitrary (say zero), and observe that (i)--(iii) hold trivially while (iv)
holds by hypothesis. Otherwise, after passing to a subsequence, choose
\(a_k^{(1)}\in \C\) such that
\[
u_{F_k}(a_k^{(1)})\ge \tfrac12\|u_{F_k}\|_{L^\infty(\C)}.
\]
By \eqref{eq:I-translation-global} and the density-translation identity
\eqref{eq:u-translation-global}, the sequence \(T_{-a_k^{(1)}}F_k\) has the
same Fock norm and the corresponding density translated; relabelling, we may
assume \(a_k^{(1)}=0\). Since \(\{F_k\}\) is bounded in \(\mathcal F^2(\C)\),
the reproducing-kernel bound \(|F_k(z)|\le \|F_k\|_{\mathcal F^2}\,
e^{\pi |z|^2/2}\) (Lemma~\ref{lem:growth}(ii)) makes \(\{F_k\}\) a normal
family of entire functions on every compact subset of \(\C\). Passing to a
subsequence and applying Montel's theorem, \(F_k\to G^{(1)}\) locally
uniformly on \(\C\), with \(G^{(1)}\in\mathcal F^2(\C)\) by Fatou's lemma. By
the choice of \(a_k^{(1)}=0\),
\[
u_{G^{(1)}}(0)=\lim_{k\to\infty}u_{F_k}(0)
\ge \tfrac12 \limsup_{k\to\infty}\|u_{F_k}\|_{L^\infty(\C)}>0,
\]
so \(G^{(1)}\) is nontrivial. Set
\[
R_k^{(1)}:=F_k-G^{(1)}.
\]

\medskip
\noindent\emph{Iteration.}
Suppose profiles \(G^{(1)},\dots,G^{(J-1)}\) and centers
\(a_k^{(1)},\dots,a_k^{(J-1)}\) have been extracted, with remainder
\(R_k^{(J-1)}\). If
\[
\limsup_{k\to\infty}\|u_{R_k^{(J-1)}}\|_{L^\infty(\C)}=0,
\]
we stop, set \(G^{(j)}=0\) and \(a_k^{(j)}=0\) for all \(j\ge J\), and
properties (i)--(iv) are easily seen to hold. Otherwise choose
\(a_k^{(J)}\in\C\) so that
\[
u_{R_k^{(J-1)}}(a_k^{(J)})
\ge \tfrac12\|u_{R_k^{(J-1)}}\|_{L^\infty(\C)}.
\]
Property (i) at this stage (proved below) forces
\(|a_k^{(J)}-a_k^{(m)}|\to\infty\) for every \(m<J\). Granted this, the
sequence \(\{T_{-a_k^{(J)}}R_k^{(J-1)}\}\) is bounded in \(\mathcal F^2\) (by
unitarity of \(T_{-a_k^{(J)}}\) together with (iii) at stage \(J-1\)), so
after passing to a further subsequence and applying Montel's theorem, it
converges locally uniformly to some \(G^{(J)}\in\mathcal F^2(\C)\). By the
choice of \(a_k^{(J)}\),
\[
u_{G^{(J)}}(0)
\ge \tfrac12 \limsup_{k\to\infty}\|u_{R_k^{(J-1)}}\|_{L^\infty(\C)}>0,
\]
so \(G^{(J)}\) is nontrivial. Set
\[
R_k^{(J)}:=R_k^{(J-1)}-T_{a_k^{(J)}}G^{(J)}.
\]
Iterating either terminates at some finite step (in which case (iv) is
immediate) or produces an infinite sequence of nontrivial profiles. By a
standard diagonal extraction we obtain a single subsequence (still denoted
\(\{F_k\}\)) for which all the partial decompositions are valid.

\medskip
\noindent\emph{Proof of (i).}
Fix \(j<\ell\), and suppose for contradiction that
\(|a_k^{(\ell)}-a_k^{(j)}|\) does not tend to infinity. After passing to a
subsequence, \(|a_k^{(\ell)}-a_k^{(j)}|\le C\) for some constant \(C\). At
the \(\ell\)-th step, by construction
\[
T_{-a_k^{(\ell)}}R_k^{(\ell-1)}\to G^{(\ell)}
\qquad\text{locally uniformly on }\C,
\]
with \(G^{(\ell)}\) nontrivial. On the other hand, property (ii) verified up
to stage \(\ell-1\) (proved below) gives
\[
T_{-a_k^{(j)}}R_k^{(\ell-1)}\to 0
\qquad\text{locally uniformly on }\C.
\]
The Weyl identity
\(T_{-a_k^{(\ell)}}=e^{i\theta_k}T_{-a_k^{(j)}}T_{a_k^{(j)}-a_k^{(\ell)}}\)
(for an appropriate real phase \(\theta_k\) coming from the Weyl co-cycle, of
modulus one) shows that
\[
T_{-a_k^{(\ell)}}R_k^{(\ell-1)}(z)
=
e^{i\theta_k}\big(T_{-a_k^{(j)}}R_k^{(\ell-1)}\big)\!
\big(z-(a_k^{(\ell)}-a_k^{(j)})\big).
\]
Since \(|a_k^{(\ell)}-a_k^{(j)}|\le C\), the points
\(z-(a_k^{(\ell)}-a_k^{(j)})\) remain in a fixed compact set as \(z\)
varies over a compact set; so the local uniform convergence of
\(T_{-a_k^{(j)}}R_k^{(\ell-1)}\) to \(0\) forces
\(T_{-a_k^{(\ell)}}R_k^{(\ell-1)}\to 0\) locally uniformly as well.
This contradicts the nontriviality of \(G^{(\ell)}\). Hence
\[
|a_k^{(j)}-a_k^{(\ell)}|\to\infty
\qquad\text{whenever }j\neq \ell.
\]

\medskip
\noindent\emph{Proof of (ii).}
Fix \(J\ge 1\) and \(1\le j\le J\). By construction, at the \(j\)-th
extraction step,
\[
T_{-a_k^{(j)}}R_k^{(j-1)}\to G^{(j)}
\qquad\text{locally uniformly on }\C.
\]
Subtracting the \(j\)-th extracted profile,
\[
T_{-a_k^{(j)}}R_k^{(j)}
=
T_{-a_k^{(j)}}R_k^{(j-1)}-G^{(j)}
\to 0
\qquad\text{locally uniformly on }\C.
\]
For \(J>j\), since
\[
R_k^{(J)}
=
R_k^{(j)}-\sum_{\ell=j+1}^J T_{a_k^{(\ell)}}G^{(\ell)},
\]
the linearity of the Weyl operators yields
\[
T_{-a_k^{(j)}}R_k^{(J)}
=
T_{-a_k^{(j)}}R_k^{(j)}
-\sum_{\ell=j+1}^J e^{i\theta_{k,\ell}}\,T_{a_k^{(\ell)}-a_k^{(j)}}G^{(\ell)}
\]
for some real phases \(\theta_{k,\ell}\) coming from the Weyl co-cycle. The
first term tends to \(0\) locally uniformly by the previous paragraph. For
each fixed \(\ell>j\), part (i) gives \(|a_k^{(\ell)}-a_k^{(j)}|\to\infty\),
and the identity \(u_{T_b G}(z)=u_G(z-b)\) (i.e.\
\eqref{eq:u-translation-global}) implies that
\(T_{a_k^{(\ell)}-a_k^{(j)}}G^{(\ell)}\to 0\) locally uniformly in modulus
(its density is \(u_{G^{(\ell)}}(z-(a_k^{(\ell)}-a_k^{(j)}))\to 0\) for
\(z\) in a compact set, since \(u_{G^{(\ell)}}\in C_0(\C)\)). Thus
\[
T_{-a_k^{(j)}}R_k^{(J)}\to 0
\qquad\text{locally uniformly on }\C,
\]
which is exactly (ii).

\medskip
\noindent\emph{Proof of (iii).}
For fixed \(J\ge 1\),
\[
F_k=\sum_{j=1}^J T_{a_k^{(j)}}G^{(j)}+R_k^{(J)}.
\]
Expanding the squared \(\mathcal F^2\)-norm and using
Lemma~\ref{lem:weyl-orthogonality-global} together with the pairwise
divergence of the translation parameters from (i),
\[
\Big\|\sum_{j=1}^J T_{a_k^{(j)}}G^{(j)}\Big\|_{\mathcal F^2}^2
=
\sum_{j=1}^J \|G^{(j)}\|_{\mathcal F^2}^2+o_k(1).
\]
Moreover, since
\[
T_{-a_k^{(m)}}R_k^{(J)}\to 0
\qquad\text{locally uniformly on }\C
\]
for every \(1\le m\le J\) by (ii), and the sequence
\(\{T_{-a_k^{(m)}}R_k^{(J)}\}_k\) is bounded in \(\mathcal F^2\), the same
weak-convergence argument as in
Lemma~\ref{lem:weyl-orthogonality-global} yields
\[
\langle T_{a_k^{(m)}}G^{(m)},R_k^{(J)}\rangle_{\mathcal F^2}
=
\langle G^{(m)},T_{-a_k^{(m)}}R_k^{(J)}\rangle_{\mathcal F^2}\to 0,
\]
for every fixed \(m\). Expanding \(\|F_k\|_{\mathcal F^2}^2\) and collecting
terms,
\[
\|F_k\|_{\mathcal F^2}^2
=
\sum_{j=1}^J \|G^{(j)}\|_{\mathcal F^2}^2
+\|R_k^{(J)}\|_{\mathcal F^2}^2
+o_k(1).
\]
In particular, since \(\|F_k\|_{\mathcal F^2}\) is bounded and the terms on
the right are nonnegative, summing the profile-norm contributions over \(j\)
yields the energy bound
\begin{equation}\label{eq:energy-bound-profiles}
\sum_{j\ge 1}\|G^{(j)}\|_{\mathcal F^2}^2
\le \limsup_{k\to\infty}\|F_k\|_{\mathcal F^2}^2.
\end{equation}
This is the key compactness ingredient that controls the iteration and
ensures (iv).

\medskip
\noindent\emph{Proof of (iv).}
Set
\[
\eta_J:=\limsup_{k\to\infty}\|u_{R_k^{(J)}}\|_{L^\infty(\C)},
\]
which is monotone nonincreasing in \(J\) by construction (since
\(R_k^{(J+1)}\) is obtained from \(R_k^{(J)}\) by extracting one further
nontrivial profile). Let \(\eta_\infty:=\lim_{J\to\infty}\eta_J\); we must
show \(\eta_\infty=0\). Suppose, for contradiction, that
\(\eta_\infty\ge \delta>0\).

By the iteration, at each stage \(J\) such that \(\eta_J\ge \delta\), the
extracted profile \(G^{(J+1)}\) satisfies \(u_{G^{(J+1)}}(0)\ge \delta/2\) by
construction. By Lemma~\ref{lem:Fock-basic}(ii),
\[
\|G^{(J+1)}\|_{\mathcal F^2}^2 \ge u_{G^{(J+1)}}(0)\ge \delta/2.
\]
If the iteration produced infinitely many such stages (which is the case
under the assumption \(\eta_\infty\ge \delta\)), then
\[
\sum_{j\ge 1}\|G^{(j)}\|_{\mathcal F^2}^2=\infty,
\]
contradicting \eqref{eq:energy-bound-profiles}. Hence \(\eta_\infty=0\),
proving (iv).
\end{proof}

We now turn to the global problem. The proof of existence consists of
combining the profile decomposition with the elementary continuity properties
of \(I_F(s)\), and using the homogeneity \(I_{cF}(s)=|c|^2 I_F(s)\) recorded
in Lemma~\ref{lem:basic-properties-global}(ii) to redistribute mass to a
single nontrivial profile.

\begin{theorem}\label{thm:existence-global-maximizer}
For every \(s>0\), there exists \(F_*\in \mathcal F^2(\C)\) with
\(\|F_*\|_{\mathcal F^2}=1\) such that
\[
I_{F_*}(s)=M(s).
\]
Equivalently, there exists a measurable set \(\Omega_*\subset \C\) with
\(|\Omega_*|=s\) such that
\[
\lambda_{1,\varphi}(\Omega_*)=M(s).
\]
\end{theorem}

\begin{proof}
Let \(\{F_k\}\subset \mathcal F^2(\C)\) be a maximizing sequence:
\[
\|F_k\|_{\mathcal F^2}=1,
\qquad
I_{F_k}(s)\to M(s).
\]
By replacing \(F_k\) with a Weyl translate if necessary, and using the
translation invariance \eqref{eq:I-translation-global} together with
\eqref{eq:u-translation-global}, we may suppose that
\[
u_{F_k}(0)=\|u_{F_k}\|_{L^\infty(\C)}.
\]
Note that \(M(s)>0\): indeed, taking \(F\equiv 1\), one has
\(F\in \mathcal F^2(\C)\) with \(\|F\|_{\mathcal F^2}^2=1\) and density
\(u_F(z)=e^{-\pi |z|^2}>0\) everywhere, so for any \(s>0\), \(I_F(s)>0\).
Consequently, \(\|u_{F_k}\|_{L^\infty(\C)}\) is bounded below by a positive
constant uniformly in \(k\); otherwise, for any \(\eta>0\) we would have
\(\|u_{F_k}\|_{L^\infty}<\eta\) eventually, hence \(I_{F_k}(s)\le \eta s\),
contradicting \(I_{F_k}(s)\to M(s)>0\).

Applying Proposition~\ref{prop:profile-global}, after passing to a
subsequence we may write
\[
F_k=\sum_{j=1}^J T_{a_k^{(j)}}G^{(j)}+R_k^{(J)}
\]
for every fixed \(J\ge 1\), with the conclusions of that proposition.

\medskip
\noindent\emph{Step 1: an upper bound for the limit.}
Fix \(J\ge 1\) and let \(\Omega_k^{(J)}\) be a measurable set of measure
\(s\) such that
\[
I_{F_k}(s)=\int_{\Omega_k^{(J)}} u_{F_k}\,dA;
\]
by Lemma~\ref{lem:basic-properties-global}(i) we may take
\(\Omega_k^{(J)}\) to be a superlevel set of \(u_{F_k}\). Let
\(\varepsilon>0\). Since \(u_{G^{(j)}}\in L^1(\C)\) for every \(j\) (with
\(\int_\C u_{G^{(j)}}\,dA=\|G^{(j)}\|_{\mathcal F^2}^2\)), choose \(R>0\) so
large that
\[
\int_{\C\setminus D_R} u_{G^{(j)}}\,dA<\varepsilon
\qquad\text{for all } 1\le j\le J,
\]
and large enough that the conclusion of
Lemma~\ref{lem:pointwise-decoupling-global}(ii) holds with this choice of
\(R\). Since the translation parameters are pairwise divergent by
Proposition~\ref{prop:profile-global}(i), for \(k\) large enough the disks
\[
D_R(a_k^{(1)}),\dots,D_R(a_k^{(J)})
\]
are pairwise disjoint. Write
\[
B_{k,j}:=D_R(a_k^{(j)}),
\qquad
E_{k,j}:=\Omega_k^{(J)}\cap B_{k,j},
\qquad
E_{k,0}:=\Omega_k^{(J)}\setminus \bigcup_{j=1}^J B_{k,j}.
\]
Set
\[
s_{k,j}:=|E_{k,j}|,
\qquad 0\le j\le J.
\]
Then, since the \(\{B_{k,j}\}_{j=1}^J\) are disjoint and
\(\Omega_k^{(J)}\) has measure \(s\),
\[
\sum_{j=0}^J s_{k,j}=s.
\]

Set
\[
S_k^{(J)}:=\sum_{j=1}^J T_{a_k^{(j)}}G^{(j)}.
\]
By Proposition~\ref{prop:profile-global}(iii), the sequence
\(\{S_k^{(J)}\}_k\) is bounded in \(\mathcal F^2(\C)\), with bound
\(\le 1+o_k(1)\). By Proposition~\ref{prop:profile-global}(iv), if \(J\) is
taken sufficiently large then
\[
\limsup_{k\to\infty}\|u_{R_k^{(J)}}\|_{L^\infty(\C)}<\varepsilon^2.
\]
Lemma~\ref{lem:density-perturbation-global} therefore gives
\[
\|u_{F_k}-u_{S_k^{(J)}}\|_{L^\infty(\C)}\to 0
\qquad (k\to\infty),
\]
and in particular, after taking \(k\) large enough,
\begin{equation}\label{eq:density-approximation-global}
\|u_{F_k}-u_{S_k^{(J)}}\|_{L^\infty(\C)}<\varepsilon.
\end{equation}

Now
\[
E_{k,0}
=
\Omega_k^{(J)}\setminus \bigcup_{j=1}^J B_{k,j}
\subset
\C\setminus \bigcup_{j=1}^J D_R(a_k^{(j)}).
\]
Lemma~\ref{lem:pointwise-decoupling-global}(i) gives, for each fixed
\(j\in\{1,\dots,J\}\),
\[
\sup_{z\in B_{k,j}}
\big|u_{S_k^{(J)}}(z)-u_{T_{a_k^{(j)}}G^{(j)}}(z)\big|\to 0
\qquad (k\to\infty),
\]
while Lemma~\ref{lem:pointwise-decoupling-global}(ii) yields, for our
choice of \(R\) and for \(k\) sufficiently large,
\[
\sup_{z\in E_{k,0}}u_{S_k^{(J)}}(z)\le \varepsilon.
\]
Combining these bounds with \eqref{eq:density-approximation-global}, we
conclude that for all sufficiently large \(k\),
\begin{enumerate}
\item[(a)] on each \(B_{k,j}\),
\[
u_{F_k}(z)\le u_{T_{a_k^{(j)}}G^{(j)}}(z)+2\varepsilon;
\]
\item[(b)] on \(E_{k,0}\),
\[
u_{F_k}(z)\le 2\varepsilon.
\]
\end{enumerate}

Therefore, decomposing the integral defining \(I_{F_k}(s)\) according to the
partition \(\Omega_k^{(J)}=E_{k,0}\sqcup\bigsqcup_{j=1}^J E_{k,j}\),
\begin{align*}
I_{F_k}(s)
&=
\int_{\Omega_k^{(J)}} u_{F_k}\,dA \\
&\le
\sum_{j=1}^J \int_{E_{k,j}} u_{T_{a_k^{(j)}}G^{(j)}}\,dA
+2\varepsilon \sum_{j=0}^J |E_{k,j}| \\
&=
\sum_{j=1}^J \int_{E_{k,j}-a_k^{(j)}} u_{G^{(j)}}\,dA
+2\varepsilon s \\
&\le
\sum_{j=1}^J I_{G^{(j)}}(s_{k,j})
+2\varepsilon s,
\end{align*}
where in the next-to-last step we used the change of variables
\(z\mapsto z-a_k^{(j)}\) together with \eqref{eq:u-translation-global}, and
in the last step the definition of \(I_{G^{(j)}}(s_{k,j})\) as the supremum
of \(\int_E u_{G^{(j)}}\,dA\) over sets of measure
\(s_{k,j}=|E_{k,j}|=|E_{k,j}-a_k^{(j)}|\).

Passing to a further subsequence if necessary, we may suppose
\[
s_{k,j}\to \sigma_j\in [0,s]
\qquad (k\to\infty)
\]
for each \(1\le j\le J\). By the continuity of \(I_{G^{(j)}}\) in the
measure parameter, given by Lemma~\ref{lem:I-continuity-global},
\[
M(s)=\lim_{k\to\infty}I_{F_k}(s)
\le
\sum_{j=1}^J I_{G^{(j)}}(\sigma_j)+2\varepsilon s.
\]
Since \(\varepsilon>0\) was arbitrary,
\begin{equation}\label{eq:M-upper-profiles-global}
M(s)\le \sum_{j=1}^J I_{G^{(j)}}(\sigma_j).
\end{equation}

\medskip
\noindent\emph{Step 2: redistribution of mass and saturation of the bound.}
We now use Lemma~\ref{lem:basic-properties-global}(ii). For each \(j\) with
\(G^{(j)}\neq 0\), set
\[
\widehat G^{(j)}:=\frac{G^{(j)}}{\|G^{(j)}\|_{\mathcal F^2}},
\qquad
\|\widehat G^{(j)}\|_{\mathcal F^2}=1,
\]
and use the homogeneity \(I_{cF}(\sigma)=|c|^2 I_F(\sigma)\) (with
\(c=\|G^{(j)}\|_{\mathcal F^2}\)) to obtain
\[
I_{G^{(j)}}(\sigma_j)
=
\|G^{(j)}\|_{\mathcal F^2}^2
\,I_{\widehat G^{(j)}}(\sigma_j).
\]
Since \(\widehat G^{(j)}\) has unit \(\mathcal F^2\)-norm, by definition of
\(M\) and the monotonicity of \(I_{\widehat G^{(j)}}\) given by
Lemma~\ref{lem:I-continuity-global} (using \(\sigma_j\le s\)),
\[
I_{\widehat G^{(j)}}(\sigma_j)\le I_{\widehat G^{(j)}}(s)\le M(s).
\]
Hence
\[
I_{G^{(j)}}(\sigma_j)\le \|G^{(j)}\|_{\mathcal F^2}^2\, M(s).
\]
For \(j\) with \(G^{(j)}=0\) the inequality is trivial. Summing over
\(1\le j\le J\) and using Proposition~\ref{prop:profile-global}(iii)
together with the normalization \(\|F_k\|_{\mathcal F^2}^2=1\),
\[
\sum_{j=1}^J I_{G^{(j)}}(\sigma_j)
\le
M(s)\sum_{j=1}^J \|G^{(j)}\|_{\mathcal F^2}^2
\le M(s).
\]
Combined with \eqref{eq:M-upper-profiles-global}, this yields
\[
M(s)\le \sum_{j=1}^J I_{G^{(j)}}(\sigma_j)\le M(s)
\]
for every \(J\). Therefore equality holds throughout for every \(J\), which
forces
\[
\sum_{j\ge 1}\|G^{(j)}\|_{\mathcal F^2}^2=1
\]
(so the residual \(R_k^{(J)}\) carries no \(\mathcal F^2\)-mass in the
limit), and moreover, for every \(j\) with \(G^{(j)}\neq 0\),
\[
I_{\widehat G^{(j)}}(\sigma_j)=M(s).
\]
By monotonicity of \(I_{\widehat G^{(j)}}\),
\(I_{\widehat G^{(j)}}(s)\ge I_{\widehat G^{(j)}}(\sigma_j)=M(s)\); since
\(M(s)\) is the supremum over unit-norm functions, we conclude
\[
I_{\widehat G^{(j)}}(s)=M(s).
\]
Thus every nonzero normalized profile is a global maximizer.

\medskip
\noindent\emph{Step 3: producing a maximizing set.}
Choose any \(j_0\) with \(G^{(j_0)}\neq 0\) (such \(j_0\) exists since
\(\sum_{j\ge 1}\|G^{(j)}\|_{\mathcal F^2}^2=1\)) and set
\(F_*:=\widehat G^{(j_0)}\), so that
\[
\|F_*\|_{\mathcal F^2}=1,
\qquad
I_{F_*}(s)=M(s).
\]
By Lemma~\ref{lem:basic-properties-global}(i), there exists a superlevel
set \(\Omega_*\) of \(u_{F_*}\) with \(|\Omega_*|=s\) and
\[
\int_{\Omega_*}u_{F_*}\,dA=I_{F_*}(s)=M(s).
\]
Then
\[
\lambda_{1,\varphi}(\Omega_*)\ge \int_{\Omega_*}u_{F_*}\,dA=M(s),
\]
while by definition of \(M(s)\) we have
\(\lambda_{1,\varphi}(\Omega_*)\le M(s)\). Hence
\[
\lambda_{1,\varphi}(\Omega_*)=M(s),
\]
and \(\Omega_*\) is a global maximizer of measure \(s\).
\end{proof}

We may now identify the global maximizers, by combining the existence theorem
above with the local-rigidity theorem of the previous section.

\begin{theorem}\label{thm:global-maximizers-are-disks}
For every \(s>0\), every global maximizer for
\[
\sup_{|\Omega|=s}\lambda_{1,\varphi}(\Omega)
\]
is, up to translation, a disk.
\end{theorem}

\begin{proof}
Let \(\Omega_*\) be a global maximizer of measure \(s\) (which exists by
Theorem~\ref{thm:existence-global-maximizer}). Then \(\Omega_*\) is, in
particular, a local maximizer of \(\lambda_{1,\varphi}\) under the same
measure constraint: every sufficiently small smooth deformation of
\(\Omega_*\) preserving the measure \(|\Omega|=s\) yields a competitor
whose value of \(\lambda_{1,\varphi}\) is at most
\(M(s)=\lambda_{1,\varphi}(\Omega_*)\). In particular, every competitor
\(E\) with \(|E|=|\Omega_*|\) and Gaussian-weighted symmetric-difference
distance \(\int_\C |\mathbf 1_{\Omega_*}-\mathbf 1_E|e^{-\pi|z|^2}\,dz<\delta_0\)
to \(\Omega_*\) (the gauge used in the local-maximizer definition of
Section~\ref{sec:local-max}) satisfies
\(\lambda_{1,\varphi}(E)\le \lambda_{1,\varphi}(\Omega_*)\), because
\(\lambda_{1,\varphi}(E)\le M(s)=\lambda_{1,\varphi}(\Omega_*)\) by the
defining property of \(M(s)\). The local-rigidity theorem proved in
the previous section therefore applies: every such local maximizer is, up
to translation, a disk. Since the problem is invariant under Weyl
translations, this characterization is sharp.
\end{proof}

As an immediate consequence, we recover the classical Nicola--Tilli theorem
\cite{Nicola-Tilli, Nicola-Tilli-2}.

\begin{corollary}
Among all measurable sets of prescribed measure, the disk maximizes Gaussian
time-frequency concentration:
\[
\lambda_{1,\varphi}(\Omega)\le \lambda_{1,\varphi}(D)
\]
whenever \(|\Omega|=|D|\) and \(D\) is a disk.
\end{corollary}

\begin{proof}
Existence of a global maximizer is given by
Theorem~\ref{thm:existence-global-maximizer}, and
Theorem~\ref{thm:global-maximizers-are-disks} shows that every such
maximizer is a disk, up to translation. Since translation by an element of
\(\C\) preserves both Lebesgue measure and the value of
\(\lambda_{1,\varphi}\), the claim follows.
\end{proof}


